\chapter{Memory subsystem}
\label{ch:mem}

\section{Baseline: reused byte-level V1 backend}
V2's first working milestones connected each Memory Manager's own
\code{prefetch\_engine.v} to the real, unmodified V1 chain
\code{int8\_memory\_access.v} $\to$ \code{memory\_interface.v} $\to$
\code{psram\_controller.v}, fetching one INT8 byte per transaction ---
exactly the contract V1's own \code{neuron\_memory.v} already used
against the same backend. This was correct and fully verified (bit-exact
end-to-end through the real PSRAM chain), but it was not the fastest
possible use of that chain.

\section{Optimization \#1 --- word-level burst reads}
\label{sec:burst}
Direct inspection of \code{int8\_memory\_access.v} shows it already
converts every 8-bit logical request into a \textbf{full 16-bit PSRAM
word access} internally (\code{mem\_addr <= addr >> 1}, one byte lane
selected via \code{lb\_n}/\code{ub\_n}) --- so a byte-at-a-time fetch was
already paying for two bytes of real PSRAM bandwidth per transaction
while using only one, and paying \code{int8\_memory\_access.v}'s own
request/wait round-trip twice for every real word instead of once.

\code{prefetch\_engine.v} (weights) and \code{activation\_cache.v}
(activations, \S\ref{sec:cache}) now talk directly to
\code{memory\_interface.v}'s own 16-bit word interface, \textbf{skipping
\code{int8\_memory\_access.v} entirely}. Both files remain frozen,
byte-for-byte unmodified V1 --- V2 simply chooses to reuse the lower
(word-level) layer of the same frozen stack instead of the byte-splitting
layer on top of it, the same precedent already set by
\code{slot\_mem\_arbiter.v} not reusing V1's own \code{mem\_arbiter.v}
verbatim.

\begin{fnnote}[Real, measured result --- single job, real PSRAM]
\begin{tabularx}{\textwidth}{C{2.2cm} C{2.4cm} C{2.4cm} C{1.6cm}}
\toprule
\rowh \thd{n\_tiles} & \thd{cycles, before} & \thd{cycles, after} & \thd{$\Delta$} \\
\midrule
1 & 166 & 84  & $-49\%$ \\
\rowa 3 & 446 & 204 & $-54\%$ \\
5 & 728 & 322 & $-56\%$ \\
\bottomrule
\end{tabularx}
Real Verilator simulation, real V1 PSRAM chain, all results still
bit-exact.
\end{fnnote}

Combined real wall-clock effect (256-neuron sustained workload, cycles
$\div$ real POST-P\&R Fmax): a \textbf{2.24--2.37$\times$} speedup across
every \code{N\_SLOTS} tested, at a negligible real Fmax cost
(unchanged at \code{N\_SLOTS}=1; $-6.2\%$ at \code{N\_SLOTS}=2; $-1.2\%$
at \code{N\_SLOTS}=4).

\begin{fnwarn}[Why not just pipeline more requests instead?]
\code{int8\_memory\_access.v}'s own \code{STATE\_IDLE} only samples a
new \code{req} once back in \code{STATE\_IDLE} after the previous
transaction's \code{mem\_ready} --- it fundamentally does not support
request pipelining. No wrapper built \emph{on top of} it can avoid
paying its round-trip cost twice per word; only bypassing it (talking to
\code{memory\_interface.v} directly) actually removes the redundancy.
This is why the fix reaches one layer lower in the stack rather than
adding queuing logic in front of the existing byte-level port.
\end{fnwarn}

\section{Optimization \#2 --- shared activation cache}
\label{sec:cache}
In the realistic dense-layer workloads this project benchmarks, many
neurons of the same layer share the \emph{exact same} activation vector.
Before this optimization, each of \code{N\_SLOTS} Memory Manager
instances re-fetched that identical vector from PSRAM independently ---
real, measured, redundant traffic on the one shared PSRAM port.
\code{activation\_cache.v} (a new, single shared instance per
\code{dataflow\_core}, not one per slot) fetches a given \code{x\_base}
vector once, tile by tile on first use, and serves every subsequent
request for the same vector directly from an on-chip buffer.

\begin{fnnote}[Real, measured result --- 256-neuron sustained workload, D-Stress]
\begin{tabularx}{\textwidth}{C{1.4cm} C{2.4cm} C{2.4cm} C{2.0cm} C{2.0cm}}
\toprule
\rowh \thd{N\_SLOTS} & \thd{cycles, burst only} & \thd{cycles, $+$cache} & \thd{Fmax, burst} & \thd{Fmax, $+$cache} \\
\midrule
1 & 348682 & 174610 & 152.44 & 131.79 \\
\rowa 2 & 307602 & 185428 & 133.58 & \textbf{87.72} \\
4 & 307346 & 184795 & 112.07 & 65.01 (\FAIL) \\
\bottomrule
\end{tabularx}
A further real 1.66--2.00$\times$ cycle reduction on top of optimization~\#1,
$\approx$4$\times$ combined vs.\ the original byte-level baseline.
\end{fnnote}

\begin{fnwarn}[Real, measured Fmax cost --- read this before raising N\_SLOTS]
The shared cache's real Fmax cost is \textbf{much steeper} than
optimization~\#1's: a single central resource with \code{N\_SLOTS}
request ports, a broadcast-capable hit-check evaluated combinationally
every cycle for every port, and a shared \code{tile\_store} array create
a genuine fan-in/routing hot spot that grows with \code{N\_SLOTS}.
\code{N\_SLOTS}=2 (recommended) still passes 80\,MHz (87.72\,MHz, margin
down from $+$67\% to $+$9.7\%); \code{N\_SLOTS}=4 \textbf{fails outright}
(65.01\,MHz). This is the central input to ch.~\ref{ch:roadmap}'s own
open work item on cache pipelining.
\end{fnwarn}

Combined real wall-clock speedup vs.\ the original byte-level baseline
(both optimizations together): \code{N\_SLOTS}=1 \textbf{3.86$\times$};
\code{N\_SLOTS}=2 \textbf{2.45$\times$} (the recommended configuration);
\code{N\_SLOTS}=4 2.29$\times$ but a real \emph{regression} versus
optimization~\#1 alone, since its own Fmax now fails 80\,MHz.

\subsection{Design notes}
Single-tag, tile-granular: a request tag mismatch invalidates the cache
and restarts filling from tile~0 for the new \code{x\_base} --- always
correct, never serves stale data, but can thrash under interleaved,
genuinely-different-\code{x\_base} concurrent traffic (not exercised by
this project's own dense-layer workloads, where sharing is real and
sustained). Requests are latched per-slot on arrival (the same
``queue, don't drop'' idiom used by the arbiter, \S\ref{sec:archcache}
of ch.~\ref{ch:arch}) and served with a broadcast ack the cycle a
matching tile becomes valid, so multiple slots pending on the same,
about-to-arrive tile are all served the same cycle.

\begin{fnnote}[Two real bugs found and fixed during implementation]
(1)~A target-bank/pending-bank race: a later handoff could queue a new
cache request (targeting a different double-buffer bank) in the same
cycle an earlier request was still awaiting its own ack, and
non-blocking-assignment ``last write wins'' semantics silently
misattributed which bank the earlier request's data landed in --- the
same bug class already found once for the weight-side
\code{pf\_target\_bank} register, fixed with the identical two-register
(pending/target) staging pattern. (2)~A zero-width Verilog replication
at \code{N\_SLOTS}=1 (\code{\{\$clog2(1)\{1'b0\}\}} $=$ \code{\{0\{...\}\}},
illegal outside a concatenation), the same class already found once in
\code{neural\_director.v} and fixed with the same width-agnostic
\code{'0} literal. Both found via real simulation, not by inspection.
\end{fnnote}

\section{Real PSRAM chain (unmodified V1)}
\code{memory\_interface.v} and \code{psram\_controller.v} are byte-for-byte
identical to V1's own copies throughout this chapter --- the real
page-mode support they already implement (fast same-page continuation,
slower cold access) is exploited more effectively by the word-level
rewrite, not changed. The real ISSI \code{IS66WVE4M16EBLL-70BLI} chip and
its board wiring are unchanged from V1 (ch.~\ref{ch:hw}).

\section{SDRAM upgrade addendum (2026-09-07) --- current, authoritative
memory architecture}
\label{sec:sdram-mem-addendum}
\begin{fnwarn}[Superseded architecture]
The PSRAM-based chain described above (\S\S\ref{sec:burst}--\ref{sec:cache})
belongs to an earlier V2 milestone. The project has since closed on a
single-external-memory architecture (real \code{decisions.log} DEC-0034):
\textbf{one SDR SDRAM device, one \code{sdram\_controller.v} instance},
serving weights, activations, AND results through
\code{sdram\_unified\_backend.v}'s two logical ports (W: 64-bit weight
read; AR: 16-bit, byte-maskable activation-read/result-write), arbitrated
2-way priority (W wins when both pending). No PSRAM, no second physical
memory device, in the current, frozen hardware path.
\end{fnwarn}

The device itself was upgraded mid-project from an 8\,MB part
(\code{AS4C4M16SA-6TIN}) to the current \textbf{AS4C32M16SB-7BIN,
64\,MB (512\,Mbit), 54-ball FBGA} --- both the row/column/bank geometry
(\code{sdram\_controller.v}'s \code{ROW\_BITS}/\code{COL\_BITS}/
\code{BANK\_BITS} parameters, now 13/10/2) and the SPI host protocol's
own address-field width (23$\to$26-bit byte address; WRITE\_JOB payload
grew 15$\to$18 bytes) changed accordingly. Full electrical/pinout data
and the complete FPGA$\leftrightarrow$SDRAM ball mapping are in
ch.~\ref{ch:hw}, \S\ref{sec:sdram-addendum} (kept in one place to avoid
two copies of the same real data).

\subsection{Real, measured clock closure}
\textbf{N\_SLOTS=4 @ 64\,MHz is the frozen production configuration}:
real \code{nextpnr-ecp5} P\&R, 8/8 tested seeds PASS (worst 66.58\,MHz,
worst WNS $+0.605$\,ns). \textbf{N\_SLOTS=8 @ 64\,MHz remains an open
item}: 5/8 seeds PASS (worst 60.12\,MHz, worst WNS $-1.009$\,ns) after
a real critical-path optimization (\code{sdram\_unified\_backend.v}'s
weight-cache hit-index encoder, rewritten from a serially-dependent
priority scan to a flat, parallel one-hot compare --- real errors.log
ERR-0029/decisions.log DEC-0040). 80\,MHz was tested with a genuinely
regenerated PLL (not merely a \code{--freq} flag) and is \textbf{not
achievable} at either processor count (0/8 seeds pass, both before and
after the ERR-0029 optimization) --- the achievable Fmax is a property
of the routed fabric, confirmed identical between the 64\,MHz- and
80\,MHz-targeted netlists. Bit-exact functional correctness (D-Stress,
256/256 neurons vs.\ golden model) is unaffected at every configuration
tested, including through this optimization.
